\begin{lemma}[$C^j$ boundary comparison with a block-diagonal representation]
    \label{lem:block_diagonal_boundary_periodic_components_from_pgvwc_comparison_boundary}
    \lean{MPSTensor.exists_blockDiagonal_boundary_chainGroundSpace_of_pgvwc_comparison_of_boundary}
    \leanok
    \uses{lem:block_diagonal_boundary_trace_decomposition_from_pgvwc_comparison,
        lem:block_diagonal_boundary_periodic_components_from_trace_decomposition_boundary}
    Assume $0<N$, $L\le N$, and $L+L_0\le N$. Suppose each block $A_j$
    is $L_0$-block-injective and satisfies
    $\sum_a A^j_aA^{j\dagger}_a=\Id$. Assume also that the simultaneous tuples
    $w\mapsto(A^1_w,\ldots,A^g_w)$ at the middle-word length $m$ span the full
    product algebra $\prod_j\MN{D_j}$. Let
    $\psi=\Gamma_N^{\oplus_j\mu_jA_j}(\bigoplus_jX_j)$ for some block-diagonal
    boundary conditions $X_j$. Suppose that whenever $\psi$ has such a
    representation, there are matrices $C^j_{i,\rho}$ such that, for every
    boundary-crossing interval $i$, every outside word $\rho$ of length $N-L$,
    and every word $\beta$ of length $i+L-N$ before the cut,
    $A^j_\beta C^j_{i,\rho} = ((\mu_j^NX_j)A^j_\beta)A^j_\rho$.
    Then there are block-diagonal boundary conditions $X_j$ such that the displayed
    representation holds and, for every $j$, the displayed component lies in
    the corresponding single-block ground space:
    \begin{align}
        \psi
        &=
        \Gamma_N^{\oplus_j\mu_jA_j}\!\left(\bigoplus_jX_j\right),
        \notag\\
        \Gamma_N^{A_j}(\mu_j^NX_j)
        &\in\mathcal G_{N,L}(A_j).
        \notag
    \end{align}
\end{lemma}

\begin{proof}
    \leanok
    \uses{lem:block_diagonal_boundary_trace_decomposition_from_pgvwc_comparison,
        lem:block_diagonal_boundary_periodic_components_from_trace_decomposition_boundary}
    Choose a block-diagonal boundary representation of $\psi$. For that
    representation,
    Lemma~\ref{lem:block_diagonal_boundary_trace_decomposition_from_pgvwc_comparison}
    gives, for every boundary-crossing interval $i$, outside word $\rho$,
    word $\beta$ before the cut, and middle word $w$,
    \begin{align}
        \sum_j\tr\!\left(A^j_\beta C^j_{i,\rho}A^j_w\right)
        &=
        \sum_j\tr\!\left(((\mu_j^NX_j)A^j_\beta)A^j_\rho A^j_w\right).
        \notag
    \end{align}
    Applying
    Lemma~\ref{lem:block_diagonal_boundary_periodic_components_from_trace_decomposition_boundary}
    to these trace comparisons gives
    $\Gamma_N^{A_j}(\mu_j^NX_j)\in\mathcal G_{N,L}(A_j)$ for every $j$.
\end{proof}

\begin{lemma}[Boundary trace comparisons give periodic-boundary single-block states]
    \label{lem:block_diagonal_boundary_periodic_components_from_trace_decomposition}
    \lean{MPSTensor.exists_blockDiagonal_boundary_chainGroundSpace_of_trace_decomposition_bnt_c1}
    \leanok
    \uses{lem:bnt_block_diagonal_open_boundary_matrices,
        lem:block_diagonal_boundary_periodic_components_from_trace_decomposition_boundary}
    With the hypotheses and notation of
    Lemma~\ref{lem:bnt_block_diagonal_open_boundary_matrices}, assume in
    addition that $L+L_0\le N$. Assume also that, for some middle-word length
    $m$, the simultaneous tuples $w\mapsto(A^1_w,\ldots,A^g_w)$, where $w$
    ranges over words of length $m$ in $\Fin{d}$, span the full product
    algebra $\prod_j\MN{D_j}$. Let
    $\psi\in\mathcal G_{N,L}(\bigoplus_j\mu_jA_j)$. Suppose that whenever
    $\psi=\Gamma_N^{\oplus_j\mu_jA_j}(\bigoplus_jX_j)$, there are matrices
    $C^j_{i,\rho}$ such that, for every boundary-crossing interval $i$, every
    outside word $\rho$ of length $N-L$, every word $\beta$ before the cut,
    and every word $w$ of length $m$,
    \begin{align}
        \sum_j\tr\!\left(A^j_\beta C^j_{i,\rho}A^j_w\right)
        &=
        \sum_j\tr\!\left(((\mu_j^NX_j)A^j_\beta)A^j_\rho A^j_w\right).
        \notag
    \end{align}
    Then there are block-diagonal boundary conditions $X_j$ such that the displayed
    representation holds and, for every $j$, the displayed component lies in
    the corresponding single-block ground space:
    \begin{align}
        \psi
        &=
        \Gamma_N^{\oplus_j\mu_jA_j}\!\left(\bigoplus_jX_j\right),
        \notag\\
        \Gamma_N^{A_j}(\mu_j^NX_j)
        &\in\mathcal G_{N,L}(A_j).
        \notag
    \end{align}
\end{lemma}

\begin{proof}
    \leanok
    \uses{lem:bnt_block_diagonal_open_boundary_matrices,
        lem:block_diagonal_boundary_periodic_components_from_trace_decomposition_boundary}
    By Lemma~\ref{lem:bnt_block_diagonal_open_boundary_matrices}, there are
    matrices $X_j$ with
    $\psi=\Gamma_N^{\oplus_j\mu_jA_j}(\bigoplus_jX_j)$.
    Lemma~\ref{lem:block_diagonal_boundary_periodic_components_from_trace_decomposition_boundary}
    applies to this representation and the assumed equality of boundary trace
    expressions. Thus,
    $\Gamma_N^{A_j}(\mu_j^NX_j)\in\mathcal G_{N,L}(A_j)$ for every $j$.
\end{proof}

\begin{lemma}[Basis-of-normal-tensors product span gives boundary-trace periodic components]
    \label{lem:block_diagonal_boundary_periodic_components_from_trace_decomposition_bnt_span}
    \lean{MPSTensor.exists_blockDiagonal_boundary_chainGroundSpace_of_trace_decomposition_bnt_c1_span}
    \leanok
    \uses{lem:unital_block_separating_product_span,
        lem:block_diagonal_boundary_periodic_components_from_trace_decomposition}
    With the hypotheses and notation of
    Lemma~\ref{lem:bnt_block_diagonal_open_boundary_matrices}, assume in
    addition that $L+L_0\le N$ and
    $L \ge (L_0+1) +(g-1)((L_0+1)+((L_0+1)+(L_0+1)))+1$.
    Put
    $m = (L_0+1) +(g-1)((L_0+1)+((L_0+1)+(L_0+1)))$.
    Let $\psi\in\mathcal G_{N,L}(\bigoplus_j\mu_jA_j)$. Suppose that whenever
    $\psi=\Gamma_N^{\oplus_j\mu_jA_j}(\bigoplus_jX_j)$, there are matrices
    $C^j_{i,\rho}$ such that, for every boundary-crossing interval $i$, every
    outside word $\rho$ of length $N-L$, every word $\beta$ before the cut,
    and every word $w$ of length $m$,
    \begin{align}
        \sum_j\tr\!\left(A^j_\beta C^j_{i,\rho}A^j_w\right)
        &=
        \sum_j\tr\!\left(((\mu_j^NX_j)A^j_\beta)A^j_\rho A^j_w\right).
        \notag
    \end{align}
    Then there are block-diagonal boundary conditions $X_j$ such that the displayed
    representation holds and, for every $j$, the displayed component lies in
    the corresponding single-block ground space:
    \begin{align}
        \psi
        &=
        \Gamma_N^{\oplus_j\mu_jA_j}\!\left(\bigoplus_jX_j\right),
        \notag\\
        \Gamma_N^{A_j}(\mu_j^NX_j)
        &\in\mathcal G_{N,L}(A_j).
        \notag
    \end{align}
\end{lemma}

\begin{proof}
    \leanok
    \uses{lem:unital_block_separating_product_span,
        lem:block_diagonal_boundary_periodic_components_from_trace_decomposition}
    Lemma~\ref{lem:unital_block_separating_product_span} gives
    $\spn\{(A^1_w,\ldots,A^g_w):|w|=m\} = \prod_j\MN{D_j}$.
    Applying
    Lemma~\ref{lem:block_diagonal_boundary_periodic_components_from_trace_decomposition}
    with this middle-word span gives the stated block-diagonal boundary
    conditions and the periodic-boundary constraint in each block.
\end{proof}

\begin{lemma}[$C^j$ boundary comparison gives periodic single-block states]
    \label{lem:block_diagonal_boundary_periodic_components_from_pgvwc_comparison}
    \lean{MPSTensor.exists_blockDiagonal_boundary_chainGroundSpace_of_pgvwc_comparison_bnt_c1}
    \leanok
    \uses{lem:bnt_block_diagonal_open_boundary_matrices,
        lem:block_diagonal_boundary_component_pgvwc_comparison_injective}
    With the hypotheses and notation of
    Lemma~\ref{lem:bnt_block_diagonal_open_boundary_matrices}, assume in
    addition that $L+L_0\le N$. Let
    $\psi\in\mathcal G_{N,L}(\bigoplus_j\mu_jA_j)$. Suppose that whenever
    $\psi=\Gamma_N^{\oplus_j\mu_jA_j}(\bigoplus_jX_j)$, there are matrices
    $C^j_{i,\rho}$ such that, for every boundary-crossing interval $i$, every
    outside word $\rho$ of length $N-L$, and every word $\beta$ before the cut,
    $A^j_\beta C^j_{i,\rho} = ((\mu_j^NX_j)A^j_\beta)A^j_\rho$.
    Then there are block-diagonal boundary conditions $X_j$ such that the displayed
    representation holds and, for every $j$, the displayed component lies in
    the corresponding single-block ground space:
    \begin{align}
        \psi
        &=
        \Gamma_N^{\oplus_j\mu_jA_j}\!\left(\bigoplus_jX_j\right),
        \notag\\
        \Gamma_N^{A_j}(\mu_j^NX_j)
        &\in\mathcal G_{N,L}(A_j).
        \notag
    \end{align}
\end{lemma}

\begin{proof}
    \leanok
    \uses{lem:bnt_block_diagonal_open_boundary_matrices,
        lem:block_diagonal_boundary_component_pgvwc_comparison_injective}
    By Lemma~\ref{lem:bnt_block_diagonal_open_boundary_matrices}, there are
    matrices $X_j$ with
    $\psi=\Gamma_N^{\oplus_j\mu_jA_j}(\bigoplus_jX_j)$.
    For these boundary conditions, the assumed $C^j$ boundary comparison
    supplies the matrices $C^j_{i,\rho}$.
    Lemma~\ref{lem:block_diagonal_boundary_component_pgvwc_comparison_injective}
    gives $\Gamma_N^{A_j}(\mu_j^NX_j)\in\mathcal G_{N,L}(A_j)$ for every $j$.
\end{proof}

\begin{theorem}[Crossing-tail spans give periodic single-block states]
    \label{thm:block_diagonal_boundary_periodic_components_from_crossing_span}
    \lean{MPSTensor.exists_blockDiagonal_boundary_chainGroundSpace_of_crossing_pgvwc_comparison_bnt_c1}
    \leanok
    \uses{lem:bnt_block_diagonal_open_boundary_matrices,
        lem:block_diagonal_boundary_periodic_components_from_crossing_boundary}
    With the hypotheses and notation of
    Lemma~\ref{lem:bnt_block_diagonal_open_boundary_matrices}, assume in
    addition that $L+L_0\le N$ and that every boundary-crossing interval
    beginning at $i$ has the simultaneous tail-word spanning property: the
    products
    \begin{align}
        w &\longmapsto (A^1_w,\ldots,A^g_w),
        \notag\\
        |w| &=N-i,
        \notag
    \end{align}
    span the product algebra $\prod_j\MN{D_j}$. Let
    $\psi\in\mathcal G_{N,L}(\bigoplus_j\mu_jA_j)$. Then there are
    block-diagonal boundary conditions $X_j$ such that the displayed
    representation holds and, for every $j$, the displayed component lies in
    the corresponding single-block ground space:
    \begin{align}
        \psi
        &=
        \Gamma_N^{\oplus_j\mu_jA_j}\!\left(\bigoplus_jX_j\right),
        \notag\\
        \Gamma_N^{A_j}(\mu_j^NX_j)
        &\in\mathcal G_{N,L}(A_j).
        \notag
    \end{align}
\end{theorem}

\begin{proof}
    \leanok
    \uses{lem:bnt_block_diagonal_open_boundary_matrices,
        lem:block_diagonal_boundary_periodic_components_from_crossing_boundary}
    Lemma~\ref{lem:bnt_block_diagonal_open_boundary_matrices} supplies
    block-diagonal boundary conditions $X_j$, and
    Lemma~\ref{lem:block_diagonal_boundary_periodic_components_from_crossing_boundary},
    applied with the assumed tail-word spans, places each
    $\Gamma_N^{A_j}(\mu_j^NX_j)$ in $\mathcal G_{N,L}(A_j)$.
\end{proof}

\begin{theorem}[Short crossing-tail spans give periodic single-block states]
    \label{thm:block_diagonal_boundary_periodic_components_from_short_crossing_span}
    \lean{MPSTensor.exists_blockDiagonal_boundary_chainGroundSpace_of_short_crossing_span_bnt_c1}
    \leanok
    \uses{thm:block_diagonal_boundary_periodic_components_from_crossing_span,
        lem:unital_block_separating_product_span}
    With the hypotheses and notation of
    Lemma~\ref{lem:bnt_block_diagonal_open_boundary_matrices}, assume in
    addition that $L+L_0\le N$. Assume also that, for every boundary-crossing
    interval beginning at $i$ whose tail length lies below the
    block-separation bound,
    $N-i < (L_0+1) +(g-1)((L_0+1)+((L_0+1)+(L_0+1)))$,
    the simultaneous products
    \begin{align}
        w &\longmapsto (A^1_w,\ldots,A^g_w),
        \notag\\
        |w| &=N-i,
        \notag
    \end{align}
    span the product algebra $\prod_j\MN{D_j}$. Let
    $\psi\in\mathcal G_{N,L}(\bigoplus_j\mu_jA_j)$. Then there are
    block-diagonal boundary conditions $X_j$ such that the displayed
    representation holds and, for every $j$, the displayed component lies in
    the corresponding single-block ground space:
    \begin{align}
        \psi
        &=
        \Gamma_N^{\oplus_j\mu_jA_j}\!\left(\bigoplus_jX_j\right),
        \notag\\
        \Gamma_N^{A_j}(\mu_j^NX_j)
        &\in\mathcal G_{N,L}(A_j).
        \notag
    \end{align}
    This derives the boundary-condition comparison of
    \cite[Theorem~12, proof lines~1446--1456]{PerezGarcia2007Matrix} at every
    crossing tail whose length reaches the displayed block-separation bound.
    % The remaining input is the span at the shorter crossing tails, recorded
    % in docs/paper-gaps/cpgsv21_block_diagonal_parent_ground_space.tex.
\end{theorem}

\begin{proof}
    \leanok
    \uses{thm:block_diagonal_boundary_periodic_components_from_crossing_span,
        lem:unital_block_separating_product_span}
    For a boundary-crossing interval beginning at $i$, either
    $N-i < (L_0+1) +(g-1)((L_0+1)+((L_0+1)+(L_0+1)))$
    or
    $N-i \ge (L_0+1) +(g-1)((L_0+1)+((L_0+1)+(L_0+1)))$.
    In the first case, the assumed short-tail span applies; in the second case,
    Lemma~\ref{lem:unital_block_separating_product_span} gives the
    simultaneous product span at length $N-i$. Hence every boundary-crossing
    interval has the tail-word spanning property, and
    Theorem~\ref{thm:block_diagonal_boundary_periodic_components_from_crossing_span}
    gives the block-diagonal boundary conditions and
    $\Gamma_N^{A_j}(\mu_j^NX_j)\in\mathcal G_{N,L}(A_j)$ for every $j$.
\end{proof}

\begin{definition}[Cyclic translation of configurations]
    \label{def:parent_cyclic_translate_configuration}
    \lean{MPSTensor.cyclicTranslateCfg}
    \leanok
    For a periodic configuration $\sigma$ on $N$ sites and $s\in\Z/N\Z$,
    define $(T_s\sigma)(k)=\sigma(k+s)$.
\end{definition}

\begin{definition}[Cyclic translation of states]
    \label{def:parent_cyclic_translate_state}
    \lean{MPSTensor.cyclicTranslateState}
    \leanok
    For a state $\psi$ on a periodic chain, define
    $(T_s\psi)(\sigma)=\psi(T_s\sigma)$.
\end{definition}

\begin{lemma}[Cyclic translation preserves the periodic chain space]
    \label{lem:parent_cyclic_translate_preserves_chain_space}
    \lean{MPSTensor.cyclicTranslateState_mem_chainGroundSpace}
    \leanok
    \uses{def:parent_cyclic_translate_configuration,
        def:parent_cyclic_translate_state}
    If $\psi\in\mathcal G_{N,L}(A)$, then
    $T_s\psi\in\mathcal G_{N,L}(A)$ for every cyclic translation $s$.
\end{lemma}

\begin{proof}
    \leanok
    A length-$L$ interval of $T_s\psi$ beginning at $i$ is the corresponding
    interval of $\psi$ beginning at $i-s$, with the outside configuration
    translated by $s$.
\end{proof}

\begin{lemma}[Translation exchanges the two words at a cut]
    \label{lem:parent_cyclic_translate_append}
    \lean{MPSTensor.cyclicTranslateCfg_fin_append}
    \leanok
    \uses{def:parent_cyclic_translate_configuration}
    If $\beta$ has length $K$ and $\alpha$ has positive length $S$, then
    $T_K(\beta\alpha)=\alpha\beta$ as configurations on the periodic chain of
    length $K+S$.
\end{lemma}

\begin{proof}
    \leanok
    This is addition modulo $K+S$, separated according to whether the site
    lies before or after the cut.
\end{proof}

\begin{lemma}[Comparison across a global cut]
    \label{lem:block_boundary_intertwiner_from_global_cut}
    \lean{MPSTensor.block_boundary_intertwines_of_cyclicTranslate_sum_groundSpaceMap_eq}
    \leanok
    \uses{def:parent_cyclic_translate_state,
        lem:parent_cyclic_translate_append,
        lem:block_matrices_equal_from_trace_pairings}
    Let $X_j,Y_j\in\MN{D_j}$, and suppose
    \begin{align}
        T_K\!\left(\sum_j\Gamma_{K+S}^{A_j}(X_j)\right)
        &=
        \sum_j\Gamma_{K+S}^{A_j}(Y_j).
        \notag
    \end{align}
    If the simultaneous block words of length $K$ span
    $\prod_j\MN{D_j}$, then, for every block $j$ and every word $\alpha$ of
    length $S$,
    $X_jA^j_\alpha = A^j_\alpha Y_j$.
\end{lemma}

\begin{proof}
    \leanok
    Evaluate the translated equality on the word $\beta\alpha$, where
    $|\beta|=K$. Cyclicity of the trace gives
    \begin{align}
        \sum_j\tr(X_jA^j_\alpha A^j_\beta)
        &=
        \sum_j\tr(A^j_\alpha Y_jA^j_\beta).
        \notag
    \end{align}
    The simultaneous spanning hypothesis separates the blocks and gives the
    asserted matrix identities.
\end{proof}

\begin{theorem}[A global change of cut gives periodic single-block states]
    \label{thm:block_diagonal_boundary_periodic_components_from_global_cut}
    \lean{MPSTensor.exists_blockDiagonal_boundary_chainGroundSpace_of_global_cut_bnt_c1_pgvwc07}
    \leanok
    \uses{lem:bnt_block_diagonal_open_boundary_matrices_source_length,
        thm:pgvwc_product_span_doubly_normalized,
        lem:parent_cyclic_translate_preserves_chain_space,
        lem:block_boundary_intertwiner_from_global_cut,
        lem:block_diagonal_boundary_component_right_boundary_matrices}
    With the hypotheses and notation of
    Lemma~\ref{lem:bnt_block_diagonal_open_boundary_matrices_source_length},
    assume in addition that $L+L_0\le N$. Thus
    \begin{align}
        g
        &\ge2,
        \notag\\
        L
        &\ge3(g-1)(L_0+1)+1,
        \notag\\
        N
        &\ge L+L_0.
        \notag
    \end{align}
    Let $\psi\in\mathcal G_{N,L}(\bigoplus_j\mu_jA_j)$. Then there are
    block-diagonal boundary conditions $X_j$ such that the displayed
    representation holds and, for every $j$, the displayed component lies in
    the corresponding single-block ground space:
    \begin{align}
        \psi
        &=
        \Gamma_N^{\oplus_j\mu_jA_j}\!\left(\bigoplus_jX_j\right),
        \notag\\
        \Gamma_N^{A_j}(\mu_j^NX_j)
        &\in\mathcal G_{N,L}(A_j).
        \notag
    \end{align}
    No simultaneous spanning assumption is imposed at the short crossing
    tails.
\end{theorem}

\begin{proof}
    \leanok
    \uses{lem:bnt_block_diagonal_open_boundary_matrices_source_length,
        thm:pgvwc_product_span_doubly_normalized,
        lem:parent_cyclic_translate_preserves_chain_space,
        lem:block_boundary_intertwiner_from_global_cut,
        lem:block_diagonal_boundary_component_right_boundary_matrices}
    Lemma~\ref{lem:bnt_block_diagonal_open_boundary_matrices_source_length}
    gives boundary matrices $X_j$. Translate the periodic state so that the
    cut moves past the final $L_0$ sites. By
    Lemma~\ref{lem:parent_cyclic_translate_preserves_chain_space}, the
    translated state remains in the periodic chain space, so applying
    Lemma~\ref{lem:bnt_block_diagonal_open_boundary_matrices_source_length}
    again yields boundary matrices $Y_j$ for the translated state. For a word
    $\alpha$ of length $L_0$, comparison of the two decompositions against
    every complementary word $w$ of length $N-L_0$ gives
    \begin{align}
        \sum_j\tr\!\left((\mu_j^NX_j)A^j_\alpha A^j_w\right)
        &=
        \sum_j\tr\!\left(A^j_\alpha(\mu_j^NY_j)A^j_w\right).
        \notag
    \end{align}
    The length bound and
    Theorem~\ref{thm:pgvwc_product_span_doubly_normalized} imply that the
    simultaneous products of length $N-L_0$ span the product algebra.
    Therefore,
    Lemma~\ref{lem:block_boundary_intertwiner_from_global_cut} gives, block by
    block,
    $(\mu_j^NX_j)A^j_\alpha = A^j_\alpha(\mu_j^NY_j)$.
    Since the words of length $L_0$ span each block matrix algebra, evaluation
    at the identity gives
    \begin{align}
        \mu_j^NX_j
        &=
        \mu_j^NY_j.
        \label{eq:ph_boundary_crossing_scaled}
    \end{align}
    Thus (\ref{eq:ph_boundary_crossing_scaled}) says that
    $\mu_j^NX_j$ commutes with every word of length $L_0$, and consequently
    with every one-site matrix. The boundary condition may then be moved
    across every interval crossing the cut, which places each block component
    in its periodic chain space. This is the change-of-cut argument of
    \cite[Theorem~12, proof lines~1276--1289 and
    1454--1456]{PerezGarcia2007Matrix}.
\end{proof}

\begin{theorem}[Periodic block decomposition from a global cut]
    \label{thm:block_diagonal_chain_space_global_cut}
    \lean{MPSTensor.chainGroundSpace_toTensorFromBlocks_eq_iSup_of_global_cut_bnt_c1_pgvwc07}
    \leanok
    \uses{thm:block_diagonal_chain_space_global_cut_independent}
    Let $g\ge2$. Let $A_j$ be irreducible left-canonical blocks with normalized
    self-overlap, pairwise inequivalent up to gauge and phase, injective at a
    common length $L_0>0$, and satisfying
    $\sum_aA^j_aA^{j\,\dagger}_a=\Id$. Let $\mu_j\ne0$ for every $j$, and assume
    \begin{align}
        L
        &\ge3(g-1)(L_0+1)+1,
        \notag\\
        N
        &\ge L+L_0.
        \notag
    \end{align}
    Then
    \begin{align}
        \mathcal G_{N,L}\!\left(\bigoplus_j\mu_jA_j\right)
        &=
        \sum_j\mathcal G_{N,L}(A_j),
        \notag
    \end{align}
    as asserted in
    \cite[Theorem~12, proof lines~1424--1456]{PerezGarcia2007Matrix}.
\end{theorem}

\begin{proof}
    \leanok
    \uses{thm:block_diagonal_chain_space_global_cut_independent}
    This is the equality component of
    Theorem~\ref{thm:block_diagonal_chain_space_global_cut_independent}.
\end{proof}

\begin{theorem}[Periodic block decomposition and block-space independence]
    \label{thm:block_diagonal_chain_space_global_cut_independent}
    \lean{MPSTensor.chainGroundSpace_toTensorFromBlocks_eq_iSup_and_iSupIndep_of_global_cut_bnt_c1_pgvwc07}
    \leanok
    \uses{thm:block_diagonal_boundary_periodic_components_from_global_cut,
        lem:block_diagonal_boundary_representation_inclusion,
        lem:block_diagonal_boundary_closing_equality,
        thm:pgvwc_block_space_independence_doubly_normalized}
    Under the same hypotheses,
    \begin{align}
        \mathcal G_{N,L}\!\left(\bigoplus_j\mu_jA_j\right)
        &=
        \sum_j\mathcal G_{N,L}(A_j).
        \notag
    \end{align}
    The open-boundary spaces $G_N(A_j)$ also form an internal sum.
\end{theorem}

\begin{proof}
    \leanok
    \uses{thm:block_diagonal_boundary_periodic_components_from_global_cut,
        lem:block_diagonal_boundary_representation_inclusion,
        lem:block_diagonal_boundary_closing_equality,
        thm:pgvwc_block_space_independence_doubly_normalized}
    Theorem~\ref{thm:block_diagonal_boundary_periodic_components_from_global_cut}
    gives the forward inclusion in the displayed equality. Every periodic
    single-block vector is a periodic vector for the block-diagonal tensor,
    which gives the reverse inclusion.
    For independence, write $\phi_j=\Gamma_N^{A^j}(X_j)$ and suppose that
    $\sum_j\phi_j=0$. Theorem~\ref{thm:pgvwc_block_space_independence_doubly_normalized} and
    nondegeneracy of the trace pairing give
    \begin{align}
        \sum_j\Gamma_N^{A^j}(X_j)=0
        &\quad\Longrightarrow\quad
        \left(\forall (M_j)_j,
            \sum_j\tr(X_jM_j)=0\right)
        \quad\Longrightarrow\quad
        \forall j,\ X_j=0.
        \notag
    \end{align}
    The last implication follows by choosing a tuple supported in one block.
    Thus every $\phi_j$ vanishes, and the open-boundary spaces form an internal
    sum.
\end{proof}

\begin{corollary}[Kernel containment from the global cut]
    \label{cor:block_diagonal_kernel_global_cut}
    \lean{MPSTensor.ker_parentHamiltonian_toTensorFromBlocks_le_bntMPSVectorSpan_of_global_cut_bnt_c1_pgvwc07}
    \leanok
    \uses{thm:block_diagonal_chain_space_global_cut,
        lem:block_diagonal_kernel_to_bnt_span_from_chain_splitting}
    Under the same hypotheses, suppose additionally that every periodic
    single-block chain space is contained in its matrix-product-vector line.
    Then
    \begin{align}
        \ker H_L^{(N)}\!\left(\bigoplus_j\mu_jA_j\right)
        &\subseteq
        \spn\{V^{(N)}(A_j):1\le j\le g\}.
        \notag
    \end{align}
\end{corollary}

\begin{proof}
    \leanok
    \uses{thm:block_diagonal_chain_space_global_cut,
        lem:block_diagonal_kernel_to_bnt_span_from_chain_splitting}
    The parent-Hamiltonian kernel is the periodic chain space. Apply the
    preceding block decomposition and then the assumed single-block
    containments.
\end{proof}

\begin{lemma}[Matrix identities give periodic single-block states]
    \label{lem:block_diagonal_boundary_periodic_components_from_complementary_identities}
    \lean{MPSTensor.exists_blockDiagonal_boundary_chainGroundSpace_of_complementary_identities_bnt_c1}
    \leanok
    \uses{lem:bnt_block_diagonal_open_boundary_matrices,
        lem:block_diagonal_boundary_component_complementary_word_identities_injective}
    With the hypotheses and notation of
    Lemma~\ref{lem:bnt_block_diagonal_open_boundary_matrices}, assume in
    addition that $L+L_0\le N$. Let
    $\psi\in\mathcal G_{N,L}(\bigoplus_j\mu_jA_j)$. Suppose that whenever
    $\psi=\Gamma_N^{\oplus_j\mu_jA_j}(\bigoplus_jX_j)$, the corresponding
    boundary conditions satisfy the block-diagonal boundary-condition
    equations: for every boundary-crossing interval $i$, every word $\beta$
    before the cut, and every outside word $\rho$ of length $N-L$,
    $\mu_j^N X_jA^j_\beta A^j_\rho = A^j_\beta E_{j,i,\rho}$.
    Then there are block-diagonal boundary conditions $X_j$ such that the displayed
    representation holds and, for every $j$, the displayed component lies in
    the corresponding single-block ground space:
    \begin{align}
        \psi
        &=
        \Gamma_N^{\oplus_j\mu_jA_j}\!\left(\bigoplus_jX_j\right),
        \notag\\
        \Gamma_N^{A_j}(\mu_j^NX_j)
        &\in\mathcal G_{N,L}(A_j).
        \notag
    \end{align}
\end{lemma}

\begin{proof}
    \leanok
    \uses{lem:bnt_block_diagonal_open_boundary_matrices,
        lem:block_diagonal_boundary_component_complementary_word_identities_injective}
    By Lemma~\ref{lem:bnt_block_diagonal_open_boundary_matrices}, there are
    matrices $X_j$ with
    $\psi=\Gamma_N^{\oplus_j\mu_jA_j}(\bigoplus_jX_j)$.
    These $X_j$ satisfy the assumed identities, so
    Lemma~\ref{lem:block_diagonal_boundary_component_complementary_word_identities_injective}
    gives $\Gamma_N^{A_j}(\mu_j^NX_j)\in\mathcal G_{N,L}(A_j)$ for every $j$.
\end{proof}
